\section{Conclusion}
\label{sec:conclusion}

% 17 Oct (v3_full): rewritten on the freeze results; one paragraph per question in their order, then the value
% statement and the future work. The join paragraph waits for the join code tree (after the 21-22 Oct rig day).

This thesis set out to test whether a centralised controller can let a UAV leave or join a team that is carrying a
cable-suspended payload in flight, and what sets the payload's transient when it does, with a centralised
load--cable planner based on \cite{tudelft} and a cable-aware tracking MPC per UAV, flown on unmodified Betaflight
UAVs over their radio links with motion capture as the only sensing.

\paragraph{RQ1: carrying with a fixed team.} On the rig four UAVs carried the 0.86\,kg ring around a circle and
through a figure-eight within 3\,cm of the path on average with the measured rod pull and the rods planned at
60$^\circ$, against about 8\,cm before; in the rig twin three UAVs carried it along five paths at peak path speeds up
to 0.3\,m/s inside frozen bars in 39 of 40 runs. In the rigid-rod twin the measured pull is also what holds the ring's
height (separated in an exploratory one-variable pair): it takes up a vertical force and an error in the typed mass at the UAVs, which leaves the planner's integral
the horizontal forces. Of the other steady errors, the two that are not forces are removed at their source and rods
typed long must be measured at the hand-over; a wrong attachment radius and the rig's pull towards $-x$ remain open.

\paragraph{RQ2: leave.} A leave is handled by resizing the planner's optimal control problem, because the attachment
geometry is a runtime parameter, and the switch is clean: no solve failed at 26 switches. On the rig three UAVs
carried the ring on through five of six commanded releases, with tilt rises of 5.3--16.9$^\circ$. In the twin the
transient grew as the support margin the leaver leaves shrank, a margin a condition checkable before the flight gives,
and in two hover cells it was smaller on the measured pull than on the planner's pull. What loses a ring is a magnet letting go, and a rule that drops the ring and lands every UAV turned the last
such loss on the rig into a controlled landing.

\paragraph{RQ3: join.} [After the join code tree and the 21--22 October rig day.] A UAV flown by another controller is
admitted on a confirmation this work observes from motion capture alone, the magnet tip moving with the plate on a
taut tether after the other controller's lift, which confirms the rig's one catch offline. On the rig the attach is
so far limited by the magnet reaching the plate, not by the admission.

\paragraph{Value added.} What is new is a centralised way to change the team under a carried payload, flown on
hardware: a detach or an attach is a switch to a pre-built solver with the attachment geometry as a runtime parameter,
so in the first tick after the event the planner solves for the load distribution of the team it actually has. On the
rig this carried the ring on through every commanded release that left an even triangle, announced and unannounced,
in hover and on a circle, an attitude transient that none of the closest works reports as a number
(Table~\ref{tab:lit_compare}). Three further results transfer beyond this rig: a geometric condition, checkable before a
flight, for whether a layout can lose a given UAV and keep the payload level, with the twin's tilt rise growing as its
margin shrinks; a map of the steady payload errors of a re-planned cascade on commercial flight controllers, with the
measured rod pull as the remedy that matters most; and a prediction ledger whose bands are frozen before the rig flights they
will score. They are of use to groups that build cooperative transport on commercial UAVs,
where the onboard force control of \cite{tudelft} is not available, and to work on fault tolerance and on teams whose
members are replaced in flight. What the result still lacks is scale, repeats and the join on the rig.

\paragraph{Future work.} The immediate steps are the attach on the rig with larger plates and longer tethers, then the
full rejoin with the collaborator's hardware; the hover and circle releases on the current defaults, scored against
their frozen bands; the magnets' hold against the angle of pull and their temperature; and the touchdown re-arm on the
rig. Beyond this thesis: a model of slack cables, so the planner can take over before the rods are taut; force sensing
at the UAVs, the onboard equivalent of the incremental inversion of \cite{tudelft}; outdoor flight without motion
capture; and larger teams, for which a pre-built solver per team size becomes a memory and build-time cost.
